\section{The \texttt{SANSQMonitor} McStas Component}
A circular detector measuring the radial average of intensity as a function
of the momentum transform in the sample.

\subsection*{Identification}
\begin{itemize}
  \item \textbf{Site:} 
  \item \textbf{Author:} S\&oslash;ren Kynde (kynde@nbi.dk) and Martin Cramer Pedersen (mcpe@nbi.dk)
  \item \textbf{Origin:} KU-Science
  \item \textbf{Date:} October 17, 2012
\end{itemize}

\subsection*{Description}
\begin{lstlisting}
A simple component simulating the scattering from a box-shaped, thin solution
of monodisperse, spherical particles.
\end{lstlisting}

\subsection*{Input parameters}
Parameters in \textbf{boldface} are required; the others are optional.

\begin{longtable}{p{0.22\textwidth}p{0.12\textwidth}p{0.46\textwidth}p{0.14\textwidth}}
\toprule
\textbf{Name} & \textbf{Unit} & \textbf{Description} & \textbf{Default} \\
\midrule
\endhead
NumberOfBins &  & Number of bins in the r (and q). & 100 \\
RFilename &  & File used for storing I(r). & "RDetector" \\
qFilename &  & File used for storing I(q). & "QDetector" \\
\textbf{RadiusDetector} & m & Radius of the detector (in the xy-plane). &  \\
\textbf{DistanceFromSample} & m & Distance from the sample to this component. &  \\
LambdaMin & AA & Max sensitivity in lambda - used to compute the highest possible value of momentum transfer, q. & 1.0 \\
Lambda0 &  & If lambda is given, the momentum transfers of all rays are computed from this value. Otherwise, instrumental effects are neglected. & 0.0 \\
restore\_neutron &  & If set to 1, the component restores the original neutron. & 0 \\
\bottomrule
\end{longtable}

\subsection*{Links}
\begin{itemize}
  \item \href{run:/home/willend/willend-McCode/mcstas-comps/contrib/SANSQMonitor.comp}{Source code} for \texttt{SANSQMonitor.comp}.
\end{itemize}
\IfFileExists{SANSQMonitor_static.tex}{\input{SANSQMonitor_static.tex}}{}